% Source: Lee, "Introduction to Riemannian Manifolds" (2nd ed.), Chapter 3 "Model Riemannian Manifolds", PDF pages 69-97 of the cited source.
% Transcribed from the cited source, with manual checks for a few pages, and organized according to leanblueprint
% conventions (semantic \label, bracketed titles, \uses dependency edges, \dcref to Lee's own numbering) below.

\section{Symmetries of Riemannian Manifolds}

For a Riemannian manifold $(M,g)$, $\operatorname{Iso}(M,g)$ acts on $M$ by composition. The manifold is \emph{homogeneous} if this action is transitive. Each isometry $\varphi$ has differentials $d\varphi_p:T_pM\to T_{\varphi(p)}M$ that are linear isometries. The isotropy subgroup $\operatorname{Iso}_p(M,g)$ fixes $p$ and acts on $T_pM$ through $I_p(\varphi)=d\varphi_p$. The manifold is \emph{isotropic at $p$} if this representation is transitive on the unit vectors of $T_pM$; it is isotropic if this holds at every point.

Let
\[
O(M)=\bigsqcup_{p\in M}\{\text{orthonormal bases of }T_pM\}.
\]
The induced action is
\begin{equation}
\varphi\cdot(b_1,\ldots,b_n)=\bigl(d\varphi_p(b_1),\ldots,d\varphi_p(b_n)\bigr).\tag{3.1}
\end{equation}
The manifold is \emph{frame-homogeneous} when this action is transitive. The converses below fail in general: the induced paraboloid is isotropic at one point only, Berger metrics are homogeneous but not isotropic, and Fubini--Study metrics are homogeneous and isotropic but not frame-homogeneous.

\begin{proposition}\label{prop:frame-homogeneous-manifolds}\dcref{ch3:3.1}
Let $(M,g)$ be a Riemannian manifold.
\begin{enumerate}
\item[(a)] If $M$ is isotropic at one point and homogeneous, then it is isotropic.
\item[(b)] If $M$ is frame-homogeneous, then it is homogeneous and isotropic.
\end{enumerate}
\end{proposition}
\begin{proof}
\sketch For (a), conjugating the isotropy action by an isometry carrying the
distinguished point to an arbitrary point transfers transitivity on unit
vectors. For (b), transitivity on orthonormal frames implies transitivity on
their base points and on their first unit vectors.
\end{proof}

\section{Euclidean Spaces}

Euclidean space is $(\mathbb R^n,\bar g)$ with the metric from (2.8). An $n$-dimensional inner product space $V$ carries the constant metric $g(v,w)=\langle v,w\rangle$ under $T_pV\cong V$; an orthonormal basis identifies it isometrically with $(\mathbb R^n,\bar g)$.

The Euclidean group is
\[
E(n)=\mathbb R^n\rtimes O(n),\qquad (b,A)(b',A')=(b+Ab',AA'),
\]
with faithful representation and action
\[
\rho(b,A)=\begin{pmatrix}A&b\\0&1\end{pmatrix}.
\]
\begin{equation}
(b,A)\cdot x=b+Ax.\tag{3.2}
\end{equation}
Translations and orthogonal linear maps preserve $\bar g$, and this action is transitive on $O(\mathbb R^n)$; Euclidean space is frame-homogeneous.

\section{Spheres}

For $R>0$, let $S^n(R)\subset\mathbb R^{n+1}$ be the radius-$R$ sphere with round metric $\bar g_R$ induced by the Euclidean metric; write $\bar g=\bar g_1$. The restriction of the $O(n+1)$ action is isometric.

\begin{proposition}
\label{prop:frame-homogeneous-sphere}
\dcref{ch3:3.2}
The group $O(n+1)$ acts transitively on $O(S^n(R))$, and thus each round sphere is frame-homogeneous.
\end{proposition}
\begin{proof}
\sketch Given $p\in S^n(R)$ and an orthonormal basis $(b_i)$ of $T_pS^n(R)$, the radial vector $p/R$ is orthogonal to all $b_i$. The matrix with columns $(b_1,\ldots,b_n,p/R)$ lies in $O(n+1)$ and sends the north-pole frame and north pole to $(b_i)$ and $p$.
\end{proof}

Metrics $g_1,g_2$ are \emph{conformally related} if $g_2=fg_1$ for positive smooth $f$. A diffeomorphism $\varphi:(M,g)\to(\widetilde M,\widetilde g)$ is conformal if $\varphi^*\widetilde g=fg$; conformal equivalence is existence of such a map. A manifold is \emph{locally conformally flat} if every point has a neighborhood conformally equivalent to an open Euclidean set.

Stereographic projection from $N=(0,R)$ is determined by
\begin{align*}
u^i&=\lambda\xi^i,\\
-R&=\lambda(\tau-R).
\tag{3.3}
\end{align*}
Thus
\begin{equation}
\sigma(\xi,\tau)=\frac{R\xi}{R-\tau},\tag{3.4}
\end{equation}
and solving back gives
\begin{equation}
\xi^i=\frac{u^i}{\lambda},\qquad \tau=R\frac{\lambda-1}{\lambda},\tag{3.5}
\end{equation}
\begin{equation}
\sigma^{-1}(u)=\left(\frac{2R^2u}{|u|^2+R^2},R\frac{|u|^2-R^2}{|u|^2+R^2}\right).\tag{3.6}
\end{equation}

\begin{proposition}[Stereographic projection is a conformal diffeomorphism between $S^n(R) \setminus \{N\}$ and $\mathbb{R}^n$]
\label{prop:stereographic-conformal-diffeomorphism}
\dcref{ch3:3.5}
The map $\sigma:S^n(R)\setminus\{N\}\to\mathbb R^n$ is a diffeomorphism and
\[
(\sigma^{-1})^*\bar g_R=\frac{4R^4}{(|u|^2+R^2)^2}\bar g.
\]
\end{proposition}
\begin{proof}
\sketch Formula (3.6) gives a smooth inverse. Differentiating its components and substituting into the Euclidean metric cancels the mixed terms and yields the displayed positive conformal factor.
\end{proof}

\begin{corollary}[Each sphere with a round metric is locally conformally flat]
\label{cor:sphere-round-metric-locally-conformally-flat}
\dcref{ch3:3.6}
\uses{prop:stereographic-conformal-diffeomorphism}
Every round sphere is locally conformally flat.
\end{corollary}
\begin{proof}
\sketch Stereographic projection handles neighborhoods away from the north pole; rotate the sphere, or project from the south pole, to handle the north pole.
\end{proof}

\section{Hyperbolic Spaces}

\begin{theorem}
\label{thm:hyperbolic-space-models}
\dcref{ch3:3.7}
Let $n>1$ and $R>0$. The following models are mutually isometric:
\begin{enumerate}
\item[(a)] The hyperboloid $\mathbb H^n(R)=\{(\xi,\tau):|\xi|^2-\tau^2=-R^2,\ \tau>0\}\subset\mathbb R^{n,1}$ with
\begin{equation}
\bar g_R=\iota^*\bar g,\qquad \bar g=\sum_{i=1}^n(d\xi^i)^2-(d\tau)^2.\tag{3.8}
\end{equation}
\item[(b)] The Beltrami--Klein ball $\mathbb K^n(R)=\{|w|<R\}$ with
\begin{equation}
\bar g_R=R^2\frac{\sum_i(dw^i)^2}{R^2-|w|^2}+R^2\frac{(\sum_iw^idw^i)^2}{(R^2-|w|^2)^2}.\tag{3.9}
\end{equation}
\item[(c)] The Poincare ball $\mathbb B^n(R)=\{|u|<R\}$ with
\[
\bar g_R=4R^4\frac{\sum_i(du^i)^2}{(R^2-|u|^2)^2}.
\]
\item[(d)] The Poincare half-space $\mathbb U^n(R)=\{(x,y):y>0\}$ with
\[
\bar g_R=R^2\frac{\sum_{i=1}^{n-1}(dx^i)^2+dy^2}{y^2}.
\]
\end{enumerate}
\end{theorem}
\begin{proof}
\sketch For $f(\xi,\tau)=|\xi|^2-\tau^2$, the Minkowski gradient is
\begin{equation}
\operatorname{grad}f=2\sum_i\xi^i\frac{\partial}{\partial\xi^i}+2\tau\frac{\partial}{\partial\tau},\tag{3.10}
\end{equation}
with squared norm $-4R^2$ on the hyperboloid; Corollary 2.71 makes its induced metric Riemannian.

Central projection $c:\mathbb H^n(R)\to\mathbb K^n(R)$ and its inverse are
\[
c(\xi,\tau)=\frac{R\xi}{\tau},\tag{3.11}
\]
\[
c^{-1}(w)=\left(\frac{Rw}{\sqrt{R^2-|w|^2}},\frac{R^2}{\sqrt{R^2-|w|^2}}\right).\tag{3.12}
\]
Differentiation shows that $c$ pulls the Klein metric back to the hyperboloid metric.

Hyperbolic stereographic projection $\pi:\mathbb H^n(R)\to\mathbb B^n(R)$ from $(0,-R)$ satisfies
\begin{align}
u^i&=\lambda\xi^i,\label{eq:3.13a}\\
R&=\lambda(\tau+R).\label{eq:3.13b}
\end{align}
Therefore
\[
\pi(\xi,\tau)=\frac{R\xi}{R+\tau},\qquad
|\pi(\xi,\tau)|^2=R^2\frac{\tau^2-R^2}{(\tau+R)^2}<R^2,
\]
\[
\pi^{-1}(u)=\left(\frac{2R^2u}{R^2-|u|^2},\frac{R(R^2+|u|^2)}{R^2-|u|^2}\right),
\]
and
\[
(\pi^{-1})^*\bar g_R=\frac{4R^4\sum_j(du^j)^2}{(R^2-|u|^2)^2}.
\]

The generalized Cayley map $\kappa:\mathbb U^n(R)\to\mathbb B^n(R)$ is, in dimension two,
\begin{equation}
\kappa(z)=iR\frac{z-iR}{z+iR},\tag{3.14}
\end{equation}
and in real coordinates, with $x\in\mathbb R^{n-1}$,
\begin{equation}
\kappa(x,y)=\left(\frac{2R^2x}{|x|^2+(y+R)^2},R\frac{|x|^2+|y|^2-R^2}{|x|^2+(y+R)^2}\right).\tag{3.15}
\end{equation}
Its inverse is
\[
\kappa^{-1}(u,v)=\left(\frac{2R^2u}{|u|^2+(v-R)^2},R\frac{R^2-|u|^2-v^2}{|u|^2+(v-R)^2}\right),
\]
and direct substitution gives $\kappa^*\bar g_R^{\mathbb B}=\bar g_R^{\mathbb U}$.
\end{proof}

We use $\mathbb H^n(R)$ for any model in the theorem and $\bar g_R$ for its corresponding metric; $R=1$ is denoted $(\mathbb H^n,\bar g)$.

\begin{corollary}[Each hyperbolic space is locally conformally flat]
\label{cor:hyperbolic-locally-conformally-flat}
\dcref{ch3:3.8}
Every hyperbolic space is locally conformally flat.
\end{corollary}
\begin{proof}
\sketch In the Poincare ball and half-space models, the displayed metrics are positive smooth factors times the Euclidean metric.
\end{proof}

Let $O(n,1)$ preserve the Minkowski form, and let $O^+(n,1)$ preserve the upper hyperboloid sheet. Its restriction acts isometrically on the hyperboloid model.

\begin{proposition}
\label{prop:lorentz-group-transitive}
\dcref{ch3:3.9}
\uses{prop:frame-homogeneous-sphere}
The group $O^+(n,1)$ acts transitively on $O(\mathbb H^n(R))$, and therefore $\mathbb H^n(R)$ is frame-homogeneous.
\end{proposition}
\begin{proof}
\sketch At $p\in\mathbb H^n(R)$, $p/R$ is Minkowski-unit and normal to the tangent space by (3.10). Appending it to an orthonormal tangent frame gives a Minkowski-orthonormal basis; the matrix with these columns lies in $O^+(n,1)$ and sends the north-pole frame to the prescribed frame.
\end{proof}
\section{Invariant Metrics on Lie Groups}

Let $G$ be a Lie group. A Riemannian metric $g$ on $G$ is \emph{left-invariant} if
$L_\varphi^*g=g$ for every $\varphi\in G$; right-invariance is defined analogously
using right translations, and \emph{bi-invariant} means both left- and
right-invariant.

\begin{lemma}
\label{lem:left-invariant-metrics-lie-group}
\dcref{ch3:3.10}
Let $\mathfrak g$ be the Lie algebra of $G$.
\begin{enumerate}[label=(\alph*)]
\item A Riemannian metric $g$ is left-invariant if and only if $g(X,Y)$ is
constant on $G$ for all $X,Y\in\mathfrak g$.
\item Under $T_eG\cong\mathfrak g$, the map $g\mapsto g_e\in\Sigma^2(T_e^*G)$
is a bijection from left-invariant Riemannian metrics on $G$ to inner products
on $\mathfrak g$.
\end{enumerate}
\end{lemma}

Consequently, every inner product on $\mathfrak g$ extends uniquely to a
left-invariant metric, and similarly for right-invariant metrics. Either kind
is homogeneous.

For $\varphi\in G$, let $C_\varphi(\psi)=\varphi\psi\varphi^{-1}$ and
$\operatorname{Ad}(\varphi)=(C_\varphi)_*\colon\mathfrak g\to\mathfrak g$.

\begin{proposition}[Bi-invariance and the Adjoint Representation]
\label{prop:bi-invariant-metric-adjoint}
\dcref{ch3:3.12}
\uses{lem:left-invariant-metrics-lie-group}
Let $g$ be a left-invariant metric on $G$, with corresponding inner product
$\langle\cdot,\cdot\rangle$ on $\mathfrak g$. Then $g$ is bi-invariant if and
only if
\[
\langle\operatorname{Ad}(\varphi)X,\operatorname{Ad}(\varphi)Y\rangle
=\langle X,Y\rangle
\]
for all $\varphi\in G$ and $X,Y\in\mathfrak g$.
\end{proposition}
\begin{proof}
\sketch For $\varphi,\psi\in G$ and $X,Y\in\mathfrak g$, left-invariance gives
$(R_{\varphi^{-1}})_*X=\operatorname{Ad}(\varphi)X$ and hence
\[
((R_{\varphi^{-1}})^*g)_\psi(X_\psi,Y_\psi)
=\langle\operatorname{Ad}(\varphi)X,
\operatorname{Ad}(\varphi)Y\rangle.
\]
Thus all right translations preserve $g$ exactly when the inner product is
$\operatorname{Ad}(G)$-invariant.
\end{proof}

For a subgroup $H\subseteq\operatorname{GL}(V)$, an inner product on $V$ is
\emph{$H$-invariant} if every element of $H$ is orthogonal for that product.

\begin{lemma}[Compactness Criterion for Invariant Inner Products]
\label{lem:invariant-inner-product-compact}
\dcref{ch3:3.13}
Let $V$ be a finite-dimensional real vector space and
$H\subseteq\operatorname{GL}(V)$ a subgroup. An $H$-invariant inner product on
$V$ exists if and only if the closure of $H$ in $\operatorname{GL}(V)$ is
compact.
\end{lemma}
\begin{proof}
\sketch If an $H$-invariant inner product exists, then $H$ lies in its compact
orthogonal group, so its closure is compact. Conversely, let $K=\overline H$ be
compact, choose an inner product $(\cdot,\cdot)_0$, and let $\mu$ be a
right-invariant volume form on $K$. Define
\[
(x,y)=\int_K(kx,ky)_0\,\mu(k).
\]
This is symmetric, bilinear, and positive definite. For $k_0\in K$, right
invariance of $\mu$ yields
\[
(k_0x,k_0y)=\int_K(kk_0x,kk_0y)_0\,\mu(k)=(x,y),
\]
so the product is $K$-invariant and therefore $H$-invariant.
\end{proof}

\begin{theorem}[Existence of Bi-invariant Metrics]
\label{thm:existence-bi-invariant-metrics}
\dcref{ch3:3.14}
A Lie group $G$ admits a bi-invariant metric if and only if
$\operatorname{Ad}(G)$ has compact closure in
$\operatorname{GL}(\mathfrak g)$.
\end{theorem}
\begin{proof}
\sketch Combine Proposition \ref{prop:bi-invariant-metric-adjoint} with Lemma
\ref{lem:invariant-inner-product-compact}.
\end{proof}

\begin{corollary}
\label{cor:compact-lie-group-bi-invariant-metric}
\dcref{ch3:3.15}
Every compact Lie group admits a bi-invariant Riemannian metric.
\end{corollary}
\begin{proof}
\sketch The continuous image $\operatorname{Ad}(G)$ of compact $G$ is compact;
apply Theorem \ref{thm:existence-bi-invariant-metrics}.
\end{proof}

If $\operatorname{Ad}(G)$ has compact closure, every adjoint orbit is bounded.
An unbounded orbit therefore obstructs a bi-invariant metric.

\begin{example}[Invariant Metrics on Lie Groups]
\label{ex:invariant-metrics-lie-groups}
\dcref{ch3:3.16}
\uses{ex:n-torus-riemannian-submanifold,ex:adjoint-representations-lie-groups,prop:berger-sphere,prop:orthogonal-group-bi-invariant-metric,thm:hyperbolic-space-models}
\begin{enumerate}[label=(\alph*)]
\item On an abelian Lie group the adjoint representation is trivial, so every
left-invariant metric is bi-invariant. This includes the Euclidean metric on
$\mathbb R^n$ and the flat metric on $\mathbb T^n$.

\item The metric induced on a Lie subgroup by a left-invariant, respectively
bi-invariant, metric is again left-invariant, respectively bi-invariant.

\item The group $\operatorname{SL}(2,\mathbb R)$ has left-invariant metrics but
no bi-invariant metric. In
$\mathfrak{sl}(2,\mathbb R)$, take
\[
X_0=\begin{pmatrix}0&1\\0&0\end{pmatrix},\qquad
A_c=\begin{pmatrix}c&0\\0&1/c\end{pmatrix}\quad(c>0).
\]
Then
\[
\operatorname{Ad}(A_c)X_0=A_cX_0A_c^{-1}
=\begin{pmatrix}0&c^2\\0&0\end{pmatrix},
\]
so the orbit is unbounded. The same obstruction applies to
$\operatorname{SL}(n,\mathbb R)$ for $n\ge2$, and the subgroup inheritance in
(b) excludes bi-invariant metrics on $\operatorname{GL}(n,\mathbb R)$ for
$n\ge2$. The group $\operatorname{GL}(1,\mathbb R)$ is abelian.

\item The map
\[
(w,z)\longmapsto
\begin{pmatrix}w&z\\-\overline z&\overline w\end{pmatrix} \tag{3.16}
\]
identifies $S^3\subseteq\mathbb C^2$ diffeomorphically with
$\operatorname{SU}(2)$; under this identification the round metric is
bi-invariant.

\item On the skew-symmetric matrices $\mathfrak o(n)$,
\[
\langle A,B\rangle=\operatorname{tr}(A^TB)
\]
is $\operatorname{Ad}(O(n))$-invariant and induces a bi-invariant metric on
$O(n)$.

\item The upper half-space $\mathbb U^n$ is a Lie group under the matrix
realization
\[
(x,y)\longleftrightarrow
\begin{pmatrix}I_{n-1}&0\\x^T&y\end{pmatrix},
\qquad x\in\mathbb R^{n-1},\quad y>0.
\]
Its hyperbolic metric $\mathfrak g_h$ is left-invariant but not right-invariant.

\item For $n\ge1$, the $(2n+1)$-dimensional Heisenberg group is
\[
H_n=\left\{
\begin{pmatrix}1&x^T&z\\0&I&y\\0&0&1\end{pmatrix}:
x,y\in\mathbb R^n,\ z\in\mathbb R
\right\}.
\]
It is nilpotent: the lower central series
$G\supseteq[G,G]\supseteq[G,[G,G]]\supseteq\cdots$ terminates at the trivial
group, where $[G_1,G_2]$ is generated by the commutators
$x_1x_2x_1^{-1}x_2^{-1}$. The group $H_n$ has left-invariant metrics but no
bi-invariant metric.

\item The group
\[
\mathrm{Sol}=\left\{
\begin{pmatrix}e^z&0&x\\0&e^{-z}&y\\0&0&1\end{pmatrix}:
x,y,z\in\mathbb R
\right\}
\]
is solvable: its derived series
$G\supseteq[G,G]\supseteq[[G,G],[G,G]]\supseteq\cdots$ terminates at the
trivial group. It is nonnilpotent and has left-invariant metrics but no
bi-invariant metric.
\end{enumerate}
\end{example}

The Lie groups admitting bi-invariant metrics are precisely the direct products
of compact groups and abelian groups.

\section{Other Homogeneous Riemannian Manifolds}

A \emph{homogeneous $G$-space} is a smooth manifold with a smooth transitive
action of a Lie group $G$.

\begin{theorem}[Existence of Invariant Metrics on Homogeneous Spaces]
\label{thm:existence-invariant-metrics-homogeneous-spaces}
\dcref{ch3:3.17}
\uses{lem:invariant-inner-product-compact}
\uses{lem:invariant-inner-product-compact}
\uses{thm:equivariant-rank-theorem}
\uses{thm:local-section-theorem}
Let $M$ be a homogeneous $G$-space, $p_0\in M$, and
$I_{p_0}\colon G_{p_0}\to\operatorname{GL}(T_{p_0}M)$ the isotropy
representation. A $G$-invariant Riemannian metric on $M$ exists if and only if
$I_{p_0}(G_{p_0})$ has compact closure in
$\operatorname{GL}(T_{p_0}M)$.
\end{theorem}
\begin{proof}
\sketch A $G$-invariant metric restricts at $p_0$ to an
$I_{p_0}(G_{p_0})$-invariant inner product, so Lemma
\ref{lem:invariant-inner-product-compact} gives necessity. Conversely, choose
such an invariant inner product $g_{p_0}$ and, for $p\in M$, choose
$\varphi\in G$ with $\varphi(p)=p_0$ and set
\[
g_p=(d\varphi_p)^*g_{p_0}.
\]
Isotropy invariance makes this independent of $\varphi$ and $G$-invariant.

For the orbit map $\pi(\psi)=\psi\cdot p_0$ and action map
$\theta_\varphi(p)=\varphi\cdot p$,
\begin{equation}
\pi\circ L_\varphi=\theta_\varphi\circ\pi. \tag{3.17}
\end{equation}
The equivariant rank theorem makes $\pi$ a submersion. The pullback
$\tau=\pi^*g$ is left-invariant, hence smooth because its coefficients in a
left-invariant frame are constant. A local section $\sigma\colon U\to G$ of
$\pi$ then gives $g|_U=\sigma^*\tau$, proving smoothness.
\end{proof}

\begin{corollary}[Compact Isotropy Gives an Invariant Metric]
\label{cor:invariant-metric-compact-isotropy}
\dcref{ch3:3.18}
\uses{thm:existence-invariant-metrics-homogeneous-spaces}
If a Lie group $G$ acts smoothly and transitively on a smooth manifold $M$ and
all isotropy groups are compact, then $M$ admits a $G$-invariant Riemannian
metric.
\end{corollary}

\section{Locally Homogeneous Riemannian Manifolds}

A Riemannian manifold $(M,g)$ is \emph{locally homogeneous} if, for every
$p,q\in M$, some isometry between neighborhoods of $p$ and $q$ maps $p$ to
$q$. It is \emph{locally frame-homogeneous} if the isometry can additionally be
chosen to map any prescribed orthonormal basis of $T_pM$ to any prescribed
orthonormal basis of $T_qM$. Homogeneity and frame homogeneity imply their local
counterparts, which also pass to open submanifolds.

\begin{proposition}[Riemannian Coverings Preserve Local Homogeneity]
\label{prop:covering-locally-homogeneous}
\dcref{ch3:3.20}
Let $(\widetilde M,\widetilde g)$ be homogeneous and let
$\pi\colon(\widetilde M,\widetilde g)\to(M,g)$ be a Riemannian covering. Then
$(M,g)$ is locally homogeneous. If $\widetilde M$ is frame-homogeneous, then
$M$ is locally frame-homogeneous.
\end{proposition}

\begin{theorem}[Uniformization of Compact Surfaces]
\label{thm:uniformization-compact-surfaces}
\dcref{ch3:3.22}
Every compact connected smooth $2$-manifold admits a locally
frame-homogeneous Riemannian metric and a Riemannian covering by
$\mathbb R^2$, the hyperbolic plane, or the round unit sphere.
\end{theorem}
\begin{proof}
\sketch By the classification of compact surfaces, it suffices to treat the
sphere, connected sums of tori, and connected sums of projective planes.
The round $S^2$ covers itself, the flat torus is covered by $\mathbb R^2$, and
$\mathbb{RP}^2$ has a two-sheeted round covering by $S^2$.

For an orientable surface of genus $n\ge2$, use the side word
$a_1b_1a_1^{-1}b_1^{-1}\cdots a_nb_na_n^{-1}b_n^{-1}$ on a regular
$4n$-gon in the Poincar\'e disk, with all interior angles $\pi/(2n)$. The
orientation-preserving isometry group contains
\begin{equation}
G=\left\{
\begin{pmatrix}\alpha&\beta\\\overline\beta&\overline\alpha\end{pmatrix}:
\alpha,\beta\in\mathbb C,\ |\alpha|^2-|\beta|^2>0
\right\}. \tag{3.18}
\end{equation}
The side-pairing transformations generate a discrete subgroup $\Gamma_n$ for
which $\mathbb B^2/\Gamma_n$ is the genus-$n$ surface.

The connected sum of two projective planes is the Klein bottle, represented by
\begin{equation}
\begin{aligned}
(x,0)&\sim(x,1) &&(0\le x\le1),\\
(0,y)&\sim(1,1-y) &&(0\le y\le1).
\end{aligned}
\tag{3.19}
\end{equation}
It is also $\mathbb R^2/\Gamma$, where
\begin{equation}
\Gamma=\left\{(b,A)\in E(2):
b=(k,\ell),\ k,\ell\in\mathbb Z,\quad
A=\begin{pmatrix}1&0\\0&(-1)^k\end{pmatrix}
\right\}. \tag{3.20}
\end{equation}
This action is free, proper, and isometric, so the quotient is flat. For a
connected sum of $n\ge3$ projective planes, use a regular hyperbolic $2n$-gon
with side word $a_1a_1a_2a_2\cdots a_na_n$ and interior angles $\pi/n$; the
side pairings lie in the full hyperbolic isometry group because the quotient is
nonorientable. Details are given in \cite{Ive92}.

Each quotient is locally frame-homogeneous by Proposition
\ref{prop:covering-locally-homogeneous}. Finally, uniqueness of smooth
structures on surfaces up to diffeomorphism allows the constructed metric to
be pulled back to the given smooth structure \cite{Mun56}.
\end{proof}

Every compact orientable $3$-manifold is a connected sum of compact manifolds,
each of which either admits a Riemannian covering by a homogeneous Riemannian
manifold or can be cut along embedded tori into pieces carrying finite-volume
locally homogeneous metrics. The simply connected homogeneous Riemannian
$3$-manifolds admitting finite-volume quotients are exactly
\[
\mathbb R^3,\quad S^3,\quad\mathbb H^3,\quad
S^2\times\mathbb R,\quad\mathbb H^2\times\mathbb R,
\]
with their standard product metrics where applicable, together with $H_1$,
$\mathrm{Sol}$, and the universal covering group of
$\operatorname{SL}(2,\mathbb R)$, each with a left-invariant metric; see
with references \cite{Thu97}, \cite{Sco83}, and
\cite[KL08,MF10,MT14]{BBBMP}.

\subsection{Symmetric Spaces}

A \emph{point reflection at $p$} is an isometry $\varphi\colon M\to M$ such
that $\varphi(p)=p$ and $d\varphi_p=-\operatorname{Id}_{T_pM}$. A connected
Riemannian manifold is a \emph{Riemannian symmetric space} if it has a point
reflection at every point. It is \emph{locally symmetric} if every point has a
neighborhood admitting a point reflection at that point. Every symmetric space
is locally symmetric, and every Riemannian symmetric space is homogeneous.

\begin{lemma}
\label{lem:homogeneous-point-reflection-symmetric}
\dcref{ch3:3.23}
A connected homogeneous Riemannian manifold with a point reflection at one
point is symmetric.
\end{lemma}
\begin{proof}
\sketch Let $\varphi$ reflect at $p$. For $q\in M$, choose an isometry $\psi$
with $\psi(p)=q$. Then $\widetilde\varphi=\psi\circ\varphi\circ\psi^{-1}$ fixes
$q$ and
\[
d\widetilde\varphi_q
=d\psi_p\circ(-\operatorname{Id}_{T_pM})\circ d(\psi^{-1})_q
=-\operatorname{Id}_{T_qM}.
\]
\end{proof}

\begin{example}[Riemannian Symmetric Spaces]
\label{ex:riemannian-symmetric-spaces}
\dcref{ch3:3.24}
\uses{lem:homogeneous-point-reflection-symmetric,ex:fubini-study-metric,prop:grassmann-stiefel-manifolds,prop:fubini-study-metric-homogeneous,prop:grassmannian-metric-homogeneous}
\begin{enumerate}[label=(\alph*)]
\item Every connected frame-homogeneous Riemannian manifold is symmetric:
frame homogeneity supplies an isometry fixing $p$ and taking an orthonormal
basis $(b_i)$ to $(-b_i)$. Hence Euclidean, spherical, and hyperbolic spaces are
symmetric.

\item If a connected Lie group $G$ has a bi-invariant metric, inversion
$\Phi(x)=x^{-1}$ satisfies $d\Phi_e=-\operatorname{Id}$ and is an isometry.
Indeed, bi-invariance and
$\Phi=R_{p^{-1}}\circ\Phi\circ L_{p^{-1}}$ give
\[
(\Phi^*g)_p
=d(L_{p^{-1}})_p^*\,d\Phi_e^*\,d(R_{p^{-1}})_e^*g_p=g_p.
\]
Thus inversion reflects at $e$, and Lemma
\ref{lem:homogeneous-point-reflection-symmetric} applies.

\item Complex projective spaces and Grassmann manifolds, with their standard
metrics, are symmetric spaces.

\item Products of symmetric spaces are symmetric for the product metric. A
symmetric space is \emph{irreducible} if it is not isometric to a product of
positive-dimensional symmetric spaces.
\end{enumerate}
\end{example}
\section{Model Pseudo-Riemannian Manifolds}

For $r,s\geq0$ and $R>0$, define the pseudosphere and pseudohyperbolic space by
\[
S^{r,s}(R)=\left\{(\xi,\tau)\in\mathbb R^{r+1,s}:
\sum_{i=1}^{r+1}(\xi^i)^2-\sum_{j=1}^{s}(\tau^j)^2=R^2\right\},
\]
\[
\mathbb H^{r,s}(R)=\left\{(\xi,\tau)\in\mathbb R^{r,s+1}:
\sum_{i=1}^{r}(\xi^i)^2-\sum_{j=1}^{s+1}(\tau^j)^2=-R^2\right\}.
\]
Equip them with the induced metrics
\[
\bar g_R^{(r,s)}=\iota^*\bar g^{(r+1,s)}
\quad\text{on }S^{r,s}(R),
\qquad
\bar g_R^{(r,s)}=\iota^*\bar g^{(r,s+1)}
\quad\text{on }\mathbb H^{r,s}(R).
\]

\begin{theorem}
\label{thm:pseudo-riemannian-signature}
\dcref{ch3:3.25}
For all $r,s\geq0$ and $R>0$, both $S^{r,s}(R)$ and
$\mathbb H^{r,s}(R)$ are pseudo-Riemannian manifolds of signature $(r,s)$.
\end{theorem}

For a pseudo-Riemannian manifold of signature $(r,s)$, call an orthonormal basis
standard-ordered when the positive vectors precede the negative vectors, so that
the dual-basis expression of the metric is
\[
(e^1)^2+\cdots+(e^r)^2-(e^{r+1})^2-\cdots-(e^{r+s})^2.
\]
Let $O(M)$ be the set of all standard-ordered orthonormal bases of tangent spaces.
The manifold is \emph{frame-homogeneous} when its isometry group acts transitively
on $O(M)$.

\begin{theorem}
\label{thm:frame-homogeneous-pseudo-riemannian}
\dcref{ch3:3.26}
All pseudo-Euclidean spaces, pseudospheres, and pseudohyperbolic spaces are
frame-homogeneous.
\end{theorem}

In signature $(n,1)$, $(S^{n,1}(R),\bar g_R^{(n,1)})$ is de Sitter space of
radius $R$, and $(\mathbb H^{n,1}(R),\bar g_R^{(n,1)})$ is anti-de Sitter space
of radius $R$.

\section{Further Results}

\begin{proposition}
\label{prop:rp-n-frame-homogeneous}
\dcref{ch3:3-2}
\uses{ex:real-projective-space-metric}
The metric on $\mathbb{RP}^n$ described in Example 2.34 is frame-homogeneous.
\end{proposition}

\begin{proposition}
\label{prop:berger-sphere}
\dcref{ch3:3-10}
\uses{ex:invariant-metrics-lie-groups}
In $\mathfrak{su}(2)$ let
\[
X=\begin{pmatrix}0&1\\-1&0\end{pmatrix},\quad
Y=\begin{pmatrix}0&i\\i&0\end{pmatrix},\quad
Z=\begin{pmatrix}i&0\\0&-i\end{pmatrix}.
\]
For $a>0$, let $g_a$ be the left-invariant metric on $\operatorname{SU}(2)$ for
which $X,Y,aZ$ are orthonormal.
\begin{enumerate}
\item[(a)] The metric $g_a$ is bi-invariant exactly when $a=1$.
\item[(b)] The map (3.16) is an isometry from
$(S^3,\bar g)$ to $(\operatorname{SU}(2),g_1)$.
\end{enumerate}
\end{proposition}

\begin{proposition}
\label{prop:orthogonal-group-bi-invariant-metric}
\dcref{ch3:3-11}
\uses{ex:invariant-metrics-lie-groups}
The formula
\[
\langle A,B\rangle=\operatorname{tr}(A^TB)
\]
defines a bi-invariant Riemannian metric on $O(n)$.
\end{proposition}

\begin{proposition}
\label{prop:fubini-study-metric-homogeneous}
\dcref{ch3:3-19}
\uses{ex:fubini-study-metric}
The Fubini--Study metric on $\mathbb{CP}^n$ is homogeneous, isotropic, and
symmetric.
\end{proposition}

\begin{proposition}
\label{prop:grassmannian-metric-homogeneous}
\dcref{ch3:3-20}
\uses{prop:grassmann-stiefel-manifolds}
The Riemannian-submersion metric on $G_k(\mathbb R^n)$ from
\cref{prop:grassmann-stiefel-manifolds} is homogeneous and symmetric.
\end{proposition}
